\par\noindent\textbf{(1)} 磁通Φ=Bℓx增加时，Lenz方向的电流I满足LĪ+RI=Bℓv；导杆受洛伦兹制动力F\_mag=−BℓI。牛顿方程m v̇=F−BℓI。方程符号使运动电动势向电流输入正功，而磁力从机械自由度抽取能量。
\par\noindent\textbf{(2)} 恒定v=v0时稳态Ī=0得I\_ss=Bℓv0/R；为维持该速度外推力F\_ss=BℓI\_ss=B²ℓ²v0/R。机械输入功F\_ssv0=(Bℓv0)²/R=I\_ss²R，等于电阻Joule热，真实装载负载可把R分为内阻与有用负载。
\par\noindent\textbf{(3)} 自由减速F=0且L=0时I=Bℓv/R，故m v̇=−B²ℓ²v/R。解v(t)=v(0)e\textasciicircum{}\{−t/τ\}，τ=mR/(B²ℓ²)，表现速度正比阻尼；若B或ℓ趋零，则τ→∞，磁制动消失。
\par\noindent\textbf{(4)} 把机械方程乘v、电路方程乘I并相加，得d[mv²/2+LI²/2]/dt=Fv−RI²。电感储能LI²/2必须纳入平衡；当L>0时电流对速度变化有记忆，机械阻尼并非瞬时比例关系，不能继续直接套L=0的单指数式。
